\chapter{Verification of Probabilistic Inference using Separation Logic}
\label{sec:future}

\emph{In this section we demonstrate how existing probabilistic separation logics (PSL) are capable of reasoning about probabilistic inference, the core of an interpreter/compiler of probabilistic programming languages (PPL).
We outline first steps and identify an overall structure of the work highlighting two future research questions.
We briefly argue why the use of PSL as opposed to reasoning directly in the meta logic on the semantics of a PPL, may be a promising direction of future research}
\fleuron{33}

\section{Motivation: what is the challenge with probabilistic inference?}

\Todo{1) Give some example of Bayesian inference, i.e. some interesting problem with a prior distribution and a posterior acquired using Bayes rule}

\Todo{Say that in some situation we have conjugate priors and obtain an analytic closed form solution}

But in general we do not. A major goal of probabilistic inference is to work with unnormalised distributions. Concretely
we may be interested in 
\begin{enumerate}
\item Taking an expectation over a distribution (SMC: Sequential Monte Carlo)
\item sampling from it or approximating a closed form solution (MCMC: Markov Chain Monte Carlo)
\item Determining a closed form equation for an approximation of the target distribution (Variational Inference)
\end{enumerate}

Here we focus exclusively on MCMC.

\subsection{Some more measure theory}
Let's make the goal of MCMC precise. We had stated that we wanted to draw samples given an unnormalised distribution $\overset{\sim}{\pi}(x)$ which corresponds to a true distribution $\pi(x)$ which cannot
be stated in closed form\footnote{``in closed form'' morally means that intractable operations like integrals are not allowed, but includes most other common mathematical operations (trigonometry, arithemtic, ...) are allowed}.

``Distribution'' could mean a few things. It can either mean a \emph{measure}, $\mu : \Sigma_A \to \mathbb{R}^+$ or a \emph{density}, $\rho : A \to \mathbb{R}^+$. Most common measures \emph{admit a density}, in which case we can
treat the two interchangeably, hence the terminology ``distribution'' refers to both. MCMC is only defined for those measures which admit a density.

Below we describe the conditions for a measure to admit a density

\Todo{Absolute continuity}

\section{Markov Chain Monte Carlo in a PPL}
Having established the problem precisely, let us describe the solution. An MCMC algorithm constructs a sequence of values, $x\_i : S$, where the state space $S$, is defined as the space over which our target distribution is defined. Intuitively it guarantees that after some
``burn in values'' at the beginning, each element of the rest of the sequence are all distributed according to the target distribution. 
The way in which these values are generated, is a using a probability kernel\footnote{Extremely convenient for us since the semantics of our PPL is also exactly a probability kernel}, known as the transition kernel, $T : S \marrow \mathcal{G} S$.
Suppose we have some starting value $x_0 : S$, then subsequent values are generated by sampling the kernel like so, using haskell like do notation in the Giry Monad.

\[
x_1 \leftarrow T(x_0); \\
x_2 \leftarrow T(x_1); \\
x_3 \leftarrow T(x_2); \\
\cdots
\]

This is exactly an unrolling of the for loop construct from Appl \ref{sec:appl}. We recall the syntax here: $for (N, M_i, i X. M_s)$ where N is the number of times to iterate, $M_i : \mathcal{G} A$ from which the initial value is sampled and
$\lambda i X. M_s : \mathbb{N} \to S \marrow \mathcal{G} S$. The last argument is a family of \emph{step kernels} indexed by the iteration index, i. In MCMC we only use a single step kernel ignoring the iteration index i. The previous value is fed
into the step kernel, and the next value is sampled. So restating MCMC as a for loop we have:

$X_N : \mathcal{G} S := for(N, dirac(0), \_i x. T(x))$

The choice of $M_i$ is arbitrary and does not affect correctness of MCMC. Above, we arbitrarily chose 0 as an example.


\subsection{Metropolis Hastings}

\[
\begin{aligned}
T &: S \marrow \mathcal{G} S \\
T x &= x' \leftarrow Q x \\
&let A = min \left(1, \frac{\text{density}_{\overset{\sim}{\pi}}(x')}{\text{density}_{\overset{\sim}{\pi}}(x)} \frac{\text{density}_{Q x}(x')}{\text{density}_{Q x'}(x)} \right) \\
&u \leftarrow \texttt{unif01} \\
&if u \leq A then (ret x') else (ret x) 
\end{aligned}
\]

\textbf{Future Research Problem 1: }

